\section{Conclusion}
\label{sec:conclusion}

We ran a small end-to-end \rlhf pipeline with dense checkpointing and applied one taxonomy, one reward model, two LLM judges, and a set of surface metrics to the same 10,560 responses.
In this setup the hypothesis we started from is only partly right.
Verbosity appears by 20\% of training, and the divergence between \rmodel score and judge preference becomes detectable around 40\%.
But verbosity plus sycophancy is 49.9\% of reward-hacking instances under the primary judge, not above 60\%; sycophancy does not emerge; verbosity plateaus instead of rising monotonically; and format changes come first, not last.
Fabrication is the failure mode most associated with responses that gain reward and lose the judge.

The main takeaway is that the behavior a reward model most visibly rewards is not necessarily the behavior that makes responses worse: here length grew everywhere, and fabricated content marked the rejected responses.

Several follow-ups would test how far this holds.
Human labels on a few hundred instances would validate or replace the judge taxonomy, which is the main blocker for H1.
A 3--7B policy trained with PPO and a higher token cap would show whether fabrication is specific to weak policies and whether the format-then-length ordering holds for other reward models.
Varying the share of opinion-bearing training prompts and the \rmodel's training data would test when sycophancy emerges.
More seeds and several run lengths would show whether onset times scale with steps, KL budget, or reward gained, and logit-based, order-averaged pairwise judging would remove position bias.
